% 中文摘要
\chapter[摘要]{摘\quad 要} 
\setlength{\parindent}{2em}

人体几何重建旨在从观测数据中恢复具有几何细节的数字人体模型，在医疗康复、混合现实、电影制作等领域具有广泛的应用前景。面向单目图像的人体几何重建方法具有低成本、易实施的优势，在消费市场展现出巨大的应用潜力。当前的方法通过结合参数化人体模型的正则化约束与神经隐式函数的细节恢复能力，实现了较为理想的重建效果。然而，这类方法在实际应用中仍然面临着一些挑战亟待解决：（1）从单目图像估计参数化人体模型参数的过程受噪声影响，现有方法对参数化人体模型噪声的鲁棒性不足，导致重建结果出现显著偏差；（2）现有方法对侧面结构的约束不足，使得重建结果在侧视角度下常出现不同程度的伪影现象。针对上述问题，本文围绕鲁棒性重建和侧面结构约束两方面展开研究，具体内容如下。

针对现有方法对参数化人体模型的噪声鲁棒性方面的不足，本文提出一种低频体素特征引导的人体几何重建方法。具体而言，本文首先采用网格体素化和体素编码提取参数化人体模型的低频体素特征，并结合符号距离场和法线图信息以构建局部联合特征，为后续的几何重建提供鲁棒的特征基础。此外，本文设计了一种空间交互隐式函数，用于建模全局低频体素特征与局部联合特征之间的空间关联性，在占用场预测阶段，该函数能有效利用全局拓扑信息，从而抑制噪声引起的局部预测偏差。实验表明所提方法在不同级别的噪声条件下均能保持较高的重建精度；在噪声级别为 0.4 时，其法线误差指标比现有最优方法降低了 62.9\%。

针对现有方法在侧面结构约束方面的不足，本文提出一种基于侧面结构约束的人体几何重建方法，在全局和局部两种尺度上优化预测的占用场。具体而言，本文设计一种全局合理性约束，通过引入判别器对预测占用场的侧面法向图几何结构进行监督，引导生成器预测更符合真实分布的占用场，提高重建结果侧面结构的合理性。此外，本文提出一种局部连贯性约束，利用程函方程约束预测的占用场梯度变化，以抑制表面的异常凹凸，提高重建结果表面的局部连贯性。实验结果表明，所提方法能重建出结构更合理、表面更连贯的人体侧面几何，在多种设置下的重建误差均优于现有主流方法。

% 近年来，随着生成式人工智能和神经渲染技术的快速发展，三维虚拟数字人脸的重建与驱动技术正朝着成本更低、流程更简化、重建效果更逼真的方向演进，使其逐步从电影、游戏等高成本制作场景向更广泛的普通用户应用普及。然而，当前的三维人脸重建与驱动技术仍面临诸多挑战：（1）在面向单视角输入的三维人脸重建任务中，现有方法依赖大量真实多视角人脸数据进行学习，以提取足够的三维人脸先验，而多视角人脸数据往往难以采集； （2）在面向单目视频的三维人脸驱动任务中，现有方法将三维人脸的各项参数耦合处理，忽视了不同参数的特异性，导致模型泛化能力差。针对上述挑战，本文的研究内容如下：

% 针对现有面向单视角输入的三维人脸重建方法依赖真实多视角数据提供先验的问题，本文提出一种几何增强的多视角数据合成方法，利用现有 3D 生成对抗网络合成多视角数据，并设计一种空间截断方法，从隐空间中采样一组更靠近空间重心区域的噪声来增强多视角数据的几何保真度，弥补现实世界中多视角数据缺失，从而解决了三维重建方法的数依赖真实多视角数据的问题。此外，该方法还设计了一种结合深度感知混合训练范式，将单视角数据作为补充，增加了训练数据的多样性，使得模型更加鲁棒，还引入深度感知判别器，通过对不同视角的合成图像进行约束，从而保证合成的三维人脸整体的立体度和自然度。实验结果表明所提方法显著优于现有方法，在同样的训练数据下，图像逼真度相比现有最好方法提升了42\%。

% 针对现有面向单目视频的三维人脸驱动方法训练参数耦合，模型泛化能力差的问题，本文提出一种 3D 高斯参数细粒度建模的人脸驱动方法，将基于 3D 高斯泼溅的三维人脸各个参数（顶点位置、颜色、旋转等）进行解耦建模。通过 FLAME（Faces Learned with an Articulated Model and Expressions）模型建立关于顶点位置的可学习高斯表情基，充分利用 FLAME 模型的表情变化先验；引入轻量多层感知网络（Multi-Layer Perceptrons）学习预测颜色、旋转参数与表情的联系，使得模型能建模人脸的微小变化。此外，本文还设计了一种多粒度监督策略，从“全局-局部”对模型进行全面约束，以保证模型在各种情况下都能精确驱动三维人脸。实验结果表明所提方法在各种设置下都优于现有方法，特别是在随机表情驱动的极端条件下，人脸相似度最高提升41\%。

\par \vspace{1em} % 空一行
\noindent \textbf{关键词：}LaTeX；模板；华南理工大学；排版规范


% 英文摘要
\chapter[Abstract]{Abstract}

Human geometry reconstruction aims to recover detailed digital human models from observed data, with wide applications in medical rehabilitation, mixed reality and film production. Monocular-image-based human geometry reconstruction methods are attractive for their low cost and ease of deployment, and thus have large practical potential in consumer scenarios. Current approaches combine the regularization of parametric human models with the detail-recovering capability of neural implicit functions to produce satisfactory reconstructions. However, they still face practical challenges: (1) the process of estimating parameters of a parametric human model from a single image is noise-sensitive, and existing methods lack sufficient robustness to model noise, which causes significant reconstruction errors; (2) constraints on lateral (side-view) structure are insufficient in current methods, causing visible artifacts in side views. This thesis addresses these problems by studying robustness-aware reconstruction and side-structure constraints; the main contributions are summarized below.

To improve robustness to noise in parametric model inputs, we propose a low-frequency voxel-feature-guided method for human geometry reconstruction. Specifically, we first voxelize the mesh and extract low-frequency voxel features via voxel encoding of the parametric model, and combine signed distance field and normal map information to construct local joint features that provide a robust basis for subsequent geometry reconstruction. In addition, we design a spatial-interaction implicit function to model the spatial correlation between global low-frequency voxel features and local joint features. During occupancy prediction, this function effectively leverages global topology information to suppress local prediction bias caused by noise. Experiments show that the proposed method maintains high reconstruction accuracy under different noise levels; at a noise level of 0.4, its normal error is reduced by 62.9\% compared with the current best method.

To address insufficient side-structure constraints, we propose a side-structure-constrained reconstruction approach that optimizes predicted occupancy fields at both global and local scales. Concretely, we introduce a global plausibility constraint by using a discriminator to supervise the geometric structure of the predicted side-view normal maps, guiding the generator to produce occupancy fields that better match the real distribution and thus improving the plausibility of side structures. Furthermore, we propose a local coherency constraint that enforces constraints on the gradient variations of the predicted occupancy field via governing differential-equation constraints, which suppresses surface anomalies and enhances local surface coherence. Experimental results demonstrate that our method reconstructs side geometries that are structurally more plausible and surfaces that are more coherent, and it outperforms mainstream methods in reconstruction error under various settings.

\par \vspace{1em}
\noindent \textbf{Keywords:} human geometry reconstruction; side-structure constraint; SCUT; Formatting
